\documentclass[12pt]{article}
\usepackage{amsmath}
\usepackage{amsthm}
\usepackage{listings}
\usepackage{mathtools}

\title{Analysis of the Bisection Method}
\author{Bernardo Brust}

\theoremstyle{plain}
\newtheorem*{theorem*}{Theorem}

\DeclarePairedDelimiter\ceil{\lceil}{\rceil}
\DeclarePairedDelimiter\floor{\lfloor}{\rfloor}

\begin{document}
\maketitle

\begin{abstract}
  Analyzing an implementation of the Bisection Method for finding roots of functions using GNU Octave~\cite{octave} in an informal manner, based on the information provided by the book ``Numerical Analysis''.~\cite{book}.
\end{abstract}

\section{Theoretical Development}\label{sec:theory}

The Bisection Method is a simple root finding technique that always converges under its input restrictions.

The Bisection Method relies on the Intermediate Value Theorem.

\begin{theorem*}
  Intermediate Value Theorem: let $f(x)$ be continuous on the interval $[a, b]$ and $k$ be a number such that $f(a) < k < f(b)$, then there is a value $c$ that satisfies $f(c) = k$.
\end{theorem*}

This is useful because of the following consequence: if we let $k = 0$ and $f(a)f(b) < 0$ (this implies that the signal of the function changes), then the Intermediate Value Theorem ensures that there exists (at least one) a $c$ such that $f(c) = 0$. The Bisection Method uses this to find the value of $c$ given the interval $[a, b]$.

\section{The Algorithm}\label{sec:algorithm}

The algorithm is informally described as follows:
\begin{enumerate}\label{alg:bisection}
        \item Check that the interval $[a, b]$ satisfies the condition $f(a)f(b) < 0$.
        \item Calculate the midpoint of the interval.
        \item If either we found a root at the midpoint or the interval is smaller than a given $\epsilon$ we can return the middle point.
        \item Otherwise, we check if we should reduce, a.k.a replace the edge with the midpoint, the left part (in the case of $f(a)f(p) < 0$) or the right part ($f(a)f(p) > 0$).
\end{enumerate}

Notice that the algorithm closely resembles the idea of Binary Searching from Computer Science, so the analysis will carry on a similar manner.

\section{Analysis of the Algorithm}\label{sec:analysis}

We can define a simple iteration rule: we have a search space $s = b - a$, after every step it reduces to $s/2$. Because we are halving every step, the number of iterations is proportional to the logarithm in base $2$ (binary logarithm as some call it) of the input, hence, in Computer Science terminology one could say: $\text{Bisection}(n) \in O(\lg(n))$. More mathematically, we can derive the following formula for the number of steps required to converge:
\begin{equation}\label{eq:upper-bound-bisection-v1}
  \frac{|b - a|}{2^n} < \epsilon
\end{equation}

Or written in terms of the binary logarithm with $n$ isolated:
\begin{equation}\label{eq:upper-bound-bisection-v2}
  n \le \ceil*{\log_2\left(\frac{b - a}{\epsilon}\right)}
\end{equation}

Notice that~\eqref{eq:upper-bound-bisection-v2} provides an upper bound, not the exact amount of iterations required, so we may get better results than the formula might indicate.

Moreover, notice that in~\ref{sec:theory} we stated that there exists \emph{at least} $1$ $c$ such that $f(c) = 0$, not exactly $1$. This means that even tho the Bisection will locate one root correctly there may be more.

\section{Implementation}\label{sec:implementation}

For the language of choice implementing numerical methods I used GNU Octave~\cite{octave}, as it's built with numerical computations in mind. One can find the source code in the GitHub repository ``numerical\_methods''~\cite{repo}, along with the sources of the analysis papers and other methods.

The implementation is almost a 1 to 1 translation of the algorithm described earlier, the only change we make is that instead of validating that the interval satisfies $f(a)f(b) < 0$, we check if the signs are different to avoid a multiplication.

\section{Results}\label{sec:results}

The final executable can receive a function, an interval and an $\epsilon$ to return the found root, along with the number of iterations taken and the expected error:
\begin{lstlisting}[caption={Sample execution of the program},label={lst:sample-run}]
> Enter f(x): x^2 - cos(x)
> Enter the interval [a, b]: [-0.2, pi / 2]
> Enter the tolerance/epsilon: 1e-12

Root found:
0.824132312301814
Iterations: 41
Estimated error:
8.052447597606260e-13
\end{lstlisting}

In the case above, if we tried our formula for the number of steps, we would have gotten $x > 40.6875$ which would point to a minimum of $41$ iterations, exactly what we got. However, we got an estimated error to $13$ decimal places, which is better than the $12$ we asked for.

\section{Conclusion and Discussion}\label{sec:conlusion}

In conclusion, the Bisection Method is a simple, stable and reliable algorithm to find solutions to equations in the form $f(x) = 0$, however, it is slow compared to other methods like Newton-Raphson, or even derivative-free ones such as the Secant Method. So it serves more as a theoretical basis to other methods rather than an actually useful one.

\bibliographystyle{ieeetr}
\bibliography{bibliography}

\end{document}
